\section{Discussion}

Les résultats expérimentaux révèlent plusieurs enseignements importants
pour la conception de systèmes RAG juridiques arabes.

\subsection{Performance des Modèles d'Embedding}

La supériorité des modèles 1024 dimensions (E5-large, BGE-M3) sur les
modèles 768 dimensions confirme l'importance de la capacité représentationnelle
pour le vocabulaire juridique arabe. Les 256 dimensions supplémentaires
permettent de coder des nuances sémantiques fines — cruciales pour
distinguer des articles proches comme le vol simple (art. 147) du vol
aggravé (art. 148) du Code Pénal marocain.

E5-large devance légèrement BGE-M3 sur toutes les métriques (MRR 0,81
contre 0,78, précision de domaine 89\% contre 85\%). Cependant, BGE-M3
offre l'avantage de supporter l'embedding multi-granularité (token, phrase,
document), ce qui pourrait être bénéfique pour les articles juridiques longs
(plusieurs centaines de mots). Le choix entre ces deux modèles dépend donc
du compromis entre performance pure (E5-large) et flexibilité (BGE-M3).

Le faible écart entre les performances cross-lingues (français MRR 0,74,
arabe formel MRR 0,81) est remarquable. Cela suggère que les modèles
d'embedding multilingues modernes ont atteint un niveau de généralisation
cross-lingue suffisant pour des applications spécialisées, réduisant le
besoin de modèles monolingues arabes dédiés.

\subsection{Apport de la Recherche Lexicale Arabe}

L'intégration du tokenizer arabe spécialisé a permis à BM25 d'atteindre
des scores non nuls (Recall@5 = 41,67\%, Hit@5 = 46,88\%), contrairement
aux résultats précédents où le dictionnaire PostgreSQL FTS échouait
totalement. BM25 est particulièrement utile pour :

\begin{enumerate}
    \item La recherche par numéro d'article (\texttt{المادة 147})
    \item Les requêtes contenant des termes techniques juridiques spécifiques
    \item Le filtrage post-recherche (exact match sur le domaine ou la loi)
\end{enumerate}

Dans RRF, BM25 contribue en apportant des résultats complémentaires que la
recherche dense ne trouve pas, ce qui explique l'amélioration du Recall@5
(59,38\% contre 56,77\%) et du Hit@5 (68,75\% contre 65,62\%). Le MRR@5
de RRF (0,533) est désormais proche de celui de Dense (0,565), confirmant
que le canal BM25 enrichit utilement la fusion.

\subsection{Fidélité des Réponses Générées}

La fidélité mesurée de 0,22 ne doit pas être interprétée comme un simple taux
d'hallucination. L'analyse qualitative des réponses générées révèle que le
LLM produit un \textbf{enrichissement contextuel légitime} non capturé par
la métrique de chevauchement lexical (token-overlap). En effet, le LLM
synthétise des informations provenant de plusieurs articles pour formuler
une conclusion juridique, cite les articles avec des formulations originales
(\textarabic{حسب المادة 513 من القانون الجنائي}), et reformule le langage
législatif en langage explicatif accessible. Ces comportements — souhaitables
pour un assistant juridique — abaissent artificiellement le score de fidélité
car ils ne correspondent à aucune sous-chaîne exacte des documents sources.

Plusieurs facteurs contribuent néanmoins à une fidélité réelle perfectible :

\begin{enumerate}
    \item \textbf{Taille du modèle} : Gemma 4 2B (2 milliards de paramètres)
    a une capacité limitée à suivre des instructions complexes de grounding
    \item \textbf{Quantification} : la quantification Q4\_K\_XL réduit
    l'empreinte de 73\% ($\sim$1,1 Go) mais peut dégrader la précision
    \item \textbf{Contexte long} : jusqu'à 5 articles juridiques par prompt
    \item \textbf{Absence de fine-tuning juridique} : le modèle n'a pas été
    spécialisé sur des données juridiques arabes
\end{enumerate}

Une évaluation par \textit{LLM-as-judge} capable de distinguer
l'enrichissement contextuel de l'hallucination \cite{gao2023faithful}
permettrait une mesure plus juste de la fidélité réelle du système.

\subsection{Exemples de Cas d'Utilisation}

Pour illustrer concrètement les forces et faiblesses du système, nous
présentons trois exemples de questions-réponses issus de l'application
réelle.

\subsubsection{Cas 1 : Requête réussie — Droit Pénal (arabe formel)}

\textbf{Question :}\par
\begin{flushleft}\textarabic{ما هي عقوبة السرقة البسيطة حسب القانون الجنائي المغربي ?}\end{flushleft}

\textbf{Articles retrouvés (RRF, top-3) :}
\begin{itemize}
    \item \textarabic{المادة 508 : عقوبة السرقة البسيطة...}
    \item \textarabic{المادة 509 : ظرف تشديد العقوبة إذا كان الفاعل خادما...}
    \item \textarabic{المادة 510 : ظرف تخفيف العقوبة لسرقة الضروريات...}
\end{itemize}

\textbf{Réponse générée (traduite de l'arabe) :}
\begin{quote}
\textit{\textbf{Résumé :} La peine pour vol simple au Maroc est l'emprisonnement
de 1 à 5 ans, conformément à l'article 508 du Code Pénal marocain.

\textbf{Analyse :} L'article 508 dispose que le vol simple est puni d'une
peine d'emprisonnement d'un à cinq ans. Si le voleur est un domestique ou
un employé de la victime (article 509), la peine est portée à 5-10 ans.
Le vol de biens de première nécessité commis par une personne dans le besoin
(article 510) peut bénéficier d'une peine atténuée.

\textbf{Verdict :} Pour un vol simple sans circonstances aggravantes, la
peine maximale est de 5 ans d'emprisonnement.

\textbf{Sources :} Article 508, 509, 510 du Code Pénal marocain.}
\end{quote}

\textbf{Analyse :} Cet exemple illustre un cas de succès complet. La question
est formulée en arabe formel avec un vocabulaire juridique standard
(\textarabic{عقوبة السرقة البسيطة}). Les trois articles pertinents sont
correctement retrouvés (Hit@3 = 100\%). Le LLM synthétise correctement
l'information en produisant un résumé, une analyse détaillée avec référence
croisée aux articles voisins, et un verdict juridique clair. La fidélité
est élevée car chaque affirmation est ancrée dans un article spécifique.

\subsubsection{Cas 2 : Requête partiellement réussie — Darija marocain}

\textbf{Question :}\par
\begin{flushleft}\textarabic{شنو هي العقوبة ديال السرقة فالمغرب؟}\end{flushleft}
(« Quelle est la peine pour le vol au Maroc ? », en Darija marocain)

\textbf{Articles retrouvés (RRF, top-3) :}
\begin{itemize}
    \item \textarabic{المادة 508 : عقوبة السرقة البسيطة...}
    \item \textarabic{المادة 241 : خيانة الأمانة...}
    \item \textarabic{المادة 147 : تعريف السرقة...}
\end{itemize}

\textbf{Réponse générée (traduite de l'arabe) :}
\begin{quote}
\textit{\textbf{Résumé :} Le vol est puni par la loi marocaine d'une peine
d'emprisonnement.

\textbf{Analyse :} L'article 147 définit le vol comme l'appropriation du
bien d'autrui sans son consentement. L'article 508 fixe la peine
d'emprisonnement pour vol simple. L'article 241 concerne la trahison de
confiance, qui est une infraction distincte du vol.

\textbf{Verdict :} La peine pour vol simple peut aller jusqu'à 5 ans
d'emprisonnement.

\textbf{Sources :} Articles 147, 508, 241 du Code Pénal marocain.}
\end{quote}

\textbf{Analyse :} La requête en Darija (\textarabic{شنو هي}) est correctement
comprise grâce à la capacité cross-lingue du modèle d'embedding. Cependant,
l'article 241 (trahison de confiance) est un faux positif : la similarité
sémantique entre \textarabic{السرقة} (vol) et \textarabic{خيانة الأمانة}
(trahison de confiance) est suffisamment élevée pour que l'article soit
retrouvé, alors qu'il concerne une infraction distincte. C'est un exemple
de \textbf{bruit sémantique} : le modèle capture la similarité thématique
(appropriation illicite) mais pas la distinction juridique précise. La
réponse reste correcte car le LLM distingue correctement les deux concepts.

\subsubsection{Cas 3 : Requête en échec — Droit Fiscal (domaine non aligné)}

\textbf{Question :}\par
\begin{flushleft}\textarabic{ما هي نسبة الضريبة على القيمة المضافة في المغرب؟}\end{flushleft}
(« Quel est le taux de TVA au Maroc ? »)

\textbf{Articles retrouvés (RRF) :} \textbf{Aucun.}
Le filtre \texttt{domain = `tax'} est appliqué dans la requête, mais les
articles fiscaux sont stockés sous \texttt{domain = `civil'} dans la base
de données. Aucun article n'est donc retrouvé.

\textbf{Réponse générée :}
\begin{quote}
\textit{Désolé, je n'ai pas trouvé d'articles de loi correspondant à votre
question dans les textes juridiques indexés. Veuillez reformuler votre
question ou vérifier le domaine juridique concerné.}
\end{quote}

\textbf{Analyse :} Cet exemple illustre une limitation critique identifiée
lors du benchmark. La question est parfaitement valide et les articles
fiscaux existent dans la base, mais une \textbf{inadéquation de taxonomie}
entre le filtre de recherche (\texttt{tax}) et l'étiquetage des documents
(\texttt{civil}) empêche toute récupération. C'est un problème
d'alignement des métadonnées, pas une limitation fondamentale du système.
La réponse du LLM est correcte sur le plan procédural (refus de répondre
en l'absence de contexte), ce qui démontre un comportement de \textit{safe
refusal} — préférable à une hallucination.

Ces trois exemples illustrent les capacités et les limites du système :
précision sur les questions en arabe formel avec vocabulaire standard,
robustesse face au Darija avec un risque de faux positifs sémantiques, et
comportement de refus sécurisé en cas d'absence de sources pertinentes.

\subsection{Passage à l'Échelle et Temps de Réponse}

Le temps total de réponse (11--20 secondes) est acceptable pour une
application de Q\&R juridique où l'utilisateur privilégie la qualité à la
rapidité. Cependant, plusieurs optimisations sont possibles :

\begin{itemize}
    \item \textbf{Embeddings asynchrones} : pré-calculer les embeddings des
    requêtes fréquentes
    \item \textbf{Cache sémantique} : mettre en cache les résultats de
    recherche pour des requêtes similaires
    \item \textbf{Quantification vectorielle} : réduire la dimension des
    embeddings (1024 $\rightarrow$ 512) via PCA
    \item \textbf{Inférence LLM locale} : servir le LLM localement plutôt
    que via tunnel Cloudflare (gain de 2--3 secondes de latence réseau)
\end{itemize}

\subsection{Limitations de l'Étude}

Cette étude présente plusieurs limitations :

\begin{enumerate}
    \item \textbf{Taille du corpus} : le benchmark d'embedding utilise
    seulement 26 articles. Un corpus plus large (centaines d'articles)
    permettrait des conclusions plus robustes.
    \item \textbf{Évaluation humaine} : la qualité des réponses générées
    a été évaluée par un seul annotateur. Une évaluation multi-juges avec
    calcul de l'accord inter-annotateur (Kappa de Cohen) serait préférable.
    \item \textbf{Version du LLM} : Gemma 4 2B était en phase de release
    candidate lors de l'expérimentation. Les versions stables futures
    pourraient offrir de meilleures performances.
    \item \textbf{Couverture juridique} : seuls six domaines de droit
    marocain sont couverts. L'extension à la totalité du corpus juridique
    marocain nécessiterait des ressources de calcul supplémentaires pour
    l'indexation vectorielle.
\end{enumerate}

\subsection{Travaux Futurs}

Plusieurs directions sont envisagées pour améliorer le système :

\begin{enumerate}
    \item \textbf{Amélioration de la fidélité} : fine-tuning de Gemma 4 2B
    sur un corpus de paires (question, articles, réponse) juridiques
    marocaines ; implémentation d'un module de vérification post-génération
    \item \textbf{Recherche multilingue avancée} : intégration de la
    recherche cross-lingue directe (requête française $\rightarrow$ articles
    arabes) via E5-large
    \item \textbf{Extension du corpus} : ingestion automatisée de
    l'ensemble du Bulletin Officiel du Maroc
    \item \textbf{Modèle de classification de domaines} : classifieur
    automatique de la spécialité juridique d'une question pour filtrer
    la recherche
    \item \textbf{Interface conversationnelle} : support du dialogue
    multi-tour avec mémoire de contexte et clarification de requêtes
    \item \textbf{Déploiement en production} : conteneurisation complète
    avec API, web, et service LLM orchestrés via Kubernetes
\end{enumerate}